\Article{15}
If $X$ is expressed by a given function of $u$ and $\lambda$, then $Y$ can be derived from it a priori. Set the integral $\int_0^u X\mathrm{d}u=T$, treating $\lambda$ as constant during integration. Clearly, if we differentiate with respect to $u$ in the same way, we obtain $\frac{\mathrm{d}(V+RT)}{\mathrm{d}u}=0$, so $V+RT$ is independent of $u$, or, equivalently, constant at all points of any one meridian. It must therefore be constant everywhere, since all meridians meet at the poles. If we set the value of $V$ at the north pole $=V^*$, we obtain
\[
T=\frac{V^*-V}{R}
\]
and hence
\[
Y=\frac{\mathrm{d}T}{\sin u\cdot\mathrm{d}\lambda}
\]
The result can also be written as
\[
Y=\frac{1}{\sin u}\int_0^u\frac{\mathrm{d}X}{\mathrm{d}\lambda}\cdot\mathrm{d}u
\]

\Article{16}
This remarkable proposition, that \textit{if the northward component of the horizontal magnetic force is given for the entire surface of the Earth, the westward (or eastward) component follows from it automatically}, holds in reverse only with a modification. Suppose $Y$ is expressed by a given function of $u$ and $\lambda$, and let $U$ denote the indefinite integral $\int\sin u\cdot Y\mathrm{d}\lambda$, with $u$ held constant during integration. Then $\frac{\mathrm{d}(V+RU)}{\mathrm{d}\lambda}=0$. or $V+RU$ is independent of $\lambda$ and is therefore, generally speaking, a function of $u$. Thus $\frac{\mathrm{d}(V+RU)}{R\mathrm{d}u}=\frac{\mathrm{d}U}{\mathrm{d}u}-X$ is also such a function. That is, the formula $\frac{\mathrm{d}U}{\mathrm{d}u}$ gives an incomplete expression for $X$, since a part containing only $u$ remains undetermined. The missing part can be supplied if, in addition to the expression for $Y$, we have the one for $X$ along any one specified meridian, or, more generally still, along any line extending from the \SourceLeaf{144}north pole to the south pole. We thus see that \textit{if we know the westward component of the horizontal magnetic force for the entire surface of the Earth and the northward component at all points of any one line from the north pole to the south pole, the latter component follows from this automatically for the whole surface of the Earth.}

\Article{17}
The preceding investigations concern only the horizontal part of the Earth's magnetic force. To include the vertical part, we must consider the problem in its full generality and treat $V$ as a function of \textit{three} variables that specify the position of an arbitrary point $O$ in space. For these we choose the distance $r$ from the center of the Earth, the angle $u$ that $r$ makes with the northern part of the Earth's axis, and the angle $\lambda$ between the plane through $r$ and the Earth's axis and a fixed meridian, counted as positive eastward.

Let the function $V$ be expanded in descending powers of $r$ in a series to which we give the following form
\[
V=\frac{RRP^0}{r}+\frac{R^3P'}{rr}+\frac{R^4P''}{r^3}+\frac{R^5P'''}{r^4}+\text{etc.}
\]
The coefficients $P^0,P',P''$ and so on are functions of $u$ and $\lambda$ here. To see how they are related to the distribution of the magnetic fluid inside the Earth, let $\mathrm{d}\mu$ be an element of this fluid and $\rho$ its distance from $O$. For $\mathrm{d}\mu$, let $r^0,u^0,\lambda^0$ have the same meaning as $r,u,\lambda$ have for $O$. Then $V=-\int\frac{\mathrm{d}\mu}{\rho}$, with the integral taken over all $\mathrm{d}\mu$. Furthermore,
\[
\rho=\BareRoot{\{rr-2rr^0(\cos u\cos u^0+\sin u\sin u^0\cos(\lambda-\lambda^0))+r^0r^0\}}
\]
and if we expand $\frac1\rho$ in the series
\[
\frac1\rho=\frac1r(T^0+T'\cdot\frac{r^0}{r}+T''\cdot\frac{r^0r^0}{rr}+\text{etc.})
\]
we obtain
\[
\begin{aligned}RRP^0&=-\int T^0\mathrm{d}\mu,\quad R^3P'=-\int T'r^0\mathrm{d}\mu,\\R^4P''&=-\int T''r^0r^0\mathrm{d}\mu\quad\text{etc.}\end{aligned}
\]
Since $T^0=1$, it follows from the fundamental assumption that the amounts of positive and negative fluid are equal in every measurable particle of the material carrying the fluid, and thus in the \textit{whole} Earth, or that $\int\mathrm{d}\mu=0$, that
\SourceLeaf{145}
\[
P^0=0
\]
so the first term of our series for $V$ drops out.

We also see that $P'$ has the form
\[
R^3P'=\alpha\cos u+\beta\sin u\cos\lambda+\gamma\sin u\sin\lambda
\]
where $\alpha=-\int r^0\cos u^0\cdot\mathrm{d}\mu$, $\beta=-\int r^0\sin u^0\cos\lambda^0\mathrm{d}\mu$,\\
and~\mbox{$\gamma=-\int r^0\sin u^0\sin\lambda^0\cdot\mathrm{d}\mu$}. Thus $-\alpha,-\beta,-\gamma$ are, by the definition given in Article 5 of the \textit{Intensitas vis magneticae}, the moments of terrestrial magnetism with respect to three mutually perpendicular axes. The first is the Earth's axis, and the second and third are the equatorial radii at longitudes $0$ and $90^0$.

We may take the general formulas for all the coefficients of the series for $\frac1\rho$ as known. For our purpose we need only observe that, with respect to $u$ and $\lambda$, the coefficients are polynomial functions of $\cos u$, $\sin u\cos\lambda$, and $\sin u\sin\lambda$, with $T''$ of the second order, $T'''$ of the third order, and so on. The same therefore holds for the coefficients $P'',P'''$ and so on.

The series for $\frac1\rho$ and for $V$ converge as long as $r$ is at least as large as $R$, or more precisely, at least as large as the radius of a sphere enclosing all the magnetic parts of the Earth.

\Article{18}
The function $V$, by its construction from $-\int\frac{\mathrm{d}\mu}{\rho}$, satisfies the following partial differential equation
\[
0=\frac{r\mathrm{d}\mathrm{d}rV}{\mathrm{d}r^2}+\frac{\mathrm{d}\mathrm{d}V}{\mathrm{d}u^2}+\operatorname{cotg}u\cdot\frac{\mathrm{d}V}{\mathrm{d}u}+\frac1{\sin u^2}\cdot\frac{\mathrm{d}\mathrm{d}V}{\mathrm{d}\lambda^2}
\]
which is a rearrangement of the familiar equation
\[
0=\frac{\mathrm{d}\mathrm{d}V}{\mathrm{d}x^2}+\frac{\mathrm{d}\mathrm{d}V}{\mathrm{d}y^2}+\frac{\mathrm{d}\mathrm{d}V}{\mathrm{d}z^2}
\]
where $x,y,z$ denote the rectangular coordinates of $O$. If we substitute the following value of $V$ into the former equation
\[
V=\frac{R^3P'}{rr}+\frac{R^4P''}{r^3}+\frac{R^5P'''}{r^4}+\text{etc.}
\]
we see that the individual coefficients $P',P'',P'''$ and so on likewise satisfy \SourceLeaf{146}partial differential equations, whose general expression is
\[
0=n(n+1)P^{\mathrm{n}}+\frac{\mathrm{d}\mathrm{d}P^{\mathrm{n}}}{\mathrm{d}u^2}+\operatorname{cotg}u\frac{\mathrm{d}P^{\mathrm{n}}}{\mathrm{d}u}+\frac1{\sin u^2}\cdot\frac{\mathrm{d}\mathrm{d}P^{\mathrm{n}}}{\mathrm{d}\lambda^2}
\]

Combining this equation with the observation in the preceding article gives the general form of $P^{\mathrm{n}}$. If we denote by $P^{\mathrm{n,m}}$ the following function of $u$
\[
\begin{aligned}
\{\cos u^{n-m}&-\frac{(n-m)(n-m-1)}{2(2n-1)}\cos u^{n-m-2}\\
&+\frac{(n-m)(n-m-1)(n-m-2)(n-m-3)}{2\cdot4(2n-1)(2n-3)}\cos u^{n-m-4}\\&-\text{etc.}\}\sin u^m
\end{aligned}
\]
then $P^{\mathrm{n}}$ has the form of a sum of $2n+1$ parts
\[
\begin{aligned}
P^{\mathrm{n}}&=g^{\mathrm{n,0}}P^{\mathrm{n,0}}+(g^{\mathrm{n,1}}\cos\lambda+h^{\mathrm{n,1}}\sin\lambda)P^{\mathrm{n,1}}\\
&\quad+(g^{\mathrm{n,2}}\cos2\lambda+h^{\mathrm{n,2}}\sin2\lambda)P^{\mathrm{n,2}}+\text{etc.}+(g^{\mathrm{n,n}}\cos n\lambda+h^{\mathrm{n,n}}\sin n\lambda)P^{\mathrm{n,n}}
\end{aligned}
\]
where $g^{\mathrm{n,0}},g^{\mathrm{n,1}},h^{\mathrm{n,1}},g^{\mathrm{n,2}}$ and so on are determinate numerical coefficients.

